\section{习题}

\subsection{基础过关题}

\begin{enumerate}

\item \textbf{（单项选择题）} 关于分页与分段存储管理，下列说法正确的是
\begin{center}
\begin{tabular}{ll}
(A) 页长由系统固定且对用户透明；段长可变，访问时需检查段内偏移是否小于段长 \\[4pt]
(B) 分段管理中段在内存中连续存放，因此不会产生外部碎片 \\[4pt]
(C) 分页的逻辑地址是一维的，因此寻址时不需要页号，只需页内偏移 \\[4pt]
(D) 段页式管理中，一次数据访问最多需要两次访存（查段表和页表）
\end{tabular}
\end{center}

\ifshowsolutions\par\noindent\textbf{解析：} 分页的页按物理单位等长划分，页长由系统决定，程序员看不到页的边界；分段的段按逻辑单位划分，长度由程序决定，因此地址变换时必须比较段内偏移与段长，越界则产生异常，故 (A) 正确。分段虽按逻辑单位划分，但各段长度不一，内存中段与段之间仍会留下无法利用的小空闲区，属于外部碎片，(B) 错误。分页地址确实是一维的，但页号仍要从逻辑地址高位提取并用于查页表，不能省略，(C) 错误。段页式地址变换依次访问段表、页表和数据，无 TLB 时需要三次访存，(D) 错误。\boxed{\textbf{答案：}(A)}\else \answerarea \fi

\item \textbf{（单项选择题）} 某分页系统采用二级页表，TLB 采用全相联映射。一次 TLB 查找耗时 $2$ ns，一次内存访问耗时 $100$ ns；TLB 缺失时需依次访问两级页表，再访问目标数据。忽略缺页与 Cache 的影响。若 TLB 命中率为 $96\%$，则平均每次访存的有效访问时间为
\begin{center}
\begin{tabular}{ll}
(A) $102$ ns & (B) $106$ ns \\[4pt]
(C) $110$ ns & (D) $202$ ns
\end{tabular}
\end{center}

\ifshowsolutions\par\noindent\textbf{解析：} TLB 命中时只需一次内存访问，耗时 $2+100=102$ ns。TLB 缺失时要读两级页表再取数据，共三次内存访问，耗时 $2+3\times100=302$ ns。于是
\[
\text{EAT}=0.96\times102+0.04\times302
=97.92+12.08=110\text{ ns}.
\]
选项 (B) 对应把缺失代价误算成“一级页表加数据”的 $202$ ns，(D) 则是全部缺失时的最坏时间。\boxed{\textbf{答案：}(C)}\else \answerarea \fi

\item \textbf{（填空题）} 某 32 位系统按字节编址，采用一级页表，页面大小为 $8$ KiB，页表项大小为 $4$ B。则页内偏移占 \underline{\qquad} 位，页号占 \underline{\qquad} 位，每个进程的页表占用 \underline{\qquad} 内存。

\ifshowsolutions\par\noindent\textbf{解析：} $8\text{ KiB}=2^{13}$ B，故页内偏移为 $13$ 位；页号位数 $=32-13=19$，即逻辑地址空间共有 $2^{19}$ 页。每页一个页表项，页表大小为
\[
2^{19}\times4\text{ B}=2^{21}\text{ B}=2\text{ MiB}.
\]
页表项数由页号位数决定，与物理内存大小无关。\boxed{\textbf{答案：}13,\ 19,\ 2\text{ MiB}}\else \answerarea \fi

\item \textbf{（单项选择题）} 关于缺页异常的处理，下列说法正确的是
\begin{center}
\begin{tabular}{ll}
(A) 缺页处理完成后，CPU 从触发缺页指令的下一条指令继续执行 \\[4pt]
(B) 一条指令的执行过程中至多发生一次缺页 \\[4pt]
(C) 缺页处理可能需要把被淘汰的脏页写回外存 \\[4pt]
(D) 缺页由外部设备异步发出，属于外部中断
\end{tabular}
\end{center}

\ifshowsolutions\par\noindent\textbf{解析：} 缺页属于故障（Fault），处理完必须重新执行触发缺页的那条指令，而不是执行下一条，(A) 错误。一条指令可能因取指跨页、多个操作数分处不同页而多次缺页，(B) 错误；缺页是 CPU 执行指令时同步产生的异常，不是外设发来的中断，(D) 错误。若没有空闲页框，需按置换算法选出牺牲页；该页修改位为 $1$ 时必须先写回外存，否则会丢失修改，(C) 正确。\boxed{\textbf{答案：}(C)}\else \answerarea \fi

\item \textbf{（单项选择题）} 关于抖动（thrashing）与工作集，下列说法正确的是
\begin{center}
\begin{tabular}{ll}
(A) 抖动由 CPU 利用率过低引起，增加并发进程数即可缓解 \\[4pt]
(B) 若所有进程工作集大小之和超过可用页框总数，系统可能因频繁换页而抖动 \\[4pt]
(C) 工作集窗口 $\Delta$ 越大，工作集中的页面数一定越少 \\[4pt]
(D) 采用全局置换可以从根本上消除抖动
\end{tabular}
\end{center}

\ifshowsolutions\par\noindent\textbf{解析：} 抖动的本质是缺页过于频繁，进程大部分时间在等待换入换出，CPU 利用率随之下降。其常见原因是分配给进程的页框数小于它当前工作集所需。增加并发进程数会进一步摊薄每个进程的页框，使抖动加剧，(A) 错误。工作集 $W(t,\Delta)$ 是区间 $[t-\Delta,t]$ 内访问过的页面集合，窗口越大，纳入的历史访问越多，集合只会扩大或不变，不会必然变小，(C) 错误。全局置换允许进程互相抢占页框，可能把别的进程正在使用的页换出，反而加剧抖动，(D) 错误。当所有进程工作集之和超过物理页框总数时，至少有一部分进程无法获得足够页框，系统进入频繁换页状态，(B) 正确。\boxed{\textbf{答案：}(B)}\else \answerarea \fi

\item \textbf{（计算题）} 某分段系统按字节编址，段内偏移从 $0$ 开始，段表如下：
\begin{center}
\begin{tabular}{|c|c|c|}
\hline
段号 & 段长（B） & 基址 \\
\hline
0 & 1024 & 2000 \\
1 & 500 & 6000 \\
2 & 2048 & 9000 \\
\hline
\end{tabular}
\end{center}
分别判断下列逻辑地址是否合法；合法的给出物理地址，非法的说明原因：
(1) $(0,800)$；(2) $(1,600)$；(3) $(2,1500)$。

\ifshowsolutions\par\noindent\textbf{解析：} 分段地址变换先查段表，再做越界检查 $0\le d<\text{段长}$，通过后物理地址 $=$ 基址 $+$ 段内偏移。
\begin{enumerate}[label=(\arabic*)]
\item 段 $0$ 的段长为 $1024$，$800<1024$，合法，物理地址 $=2000+800=2800$。
\item 段 $1$ 的段长为 $500$，$600\ge500$，偏移越界，产生越界异常，不是合法地址。
\item 段 $2$ 的段长为 $2048$，$1500<2048$，合法，物理地址 $=9000+1500=10500$。
\end{enumerate}
可见分段的偏移是与基址相加，而不是像分页那样直接拼接在页框号之后。\par\noindent\boxed{\textbf{答案：}2800;\ \text{越界};\ 10500}\else \answerarea \fi

\end{enumerate}

\subsection{真题风格与综合训练}

\begin{enumerate}[start=7]

\item \textbf{（单项选择题）} 某 32 位系统按字节编址，页面大小为 $4$ KiB，页表项大小为 $4$ B，采用二级页表，逻辑地址按“一级页号 $10$ 位、二级页号 $10$ 位、页内偏移 $12$ 位”划分。某进程只使用了两个虚页，其虚页号分别为 \texttt{0x00001} 和 \texttt{0x3FFFF}，这两个虚页的一级页号不同，且每个用到的一级页号下只建立一张二级页表。该进程的页目录与全部二级页表共占用多少内存？
\begin{center}
\begin{tabular}{ll}
(A) $4$ KiB & (B) $8$ KiB \\[4pt]
(C) $12$ KiB & (D) $4$ MiB
\end{tabular}
\end{center}

\ifshowsolutions\par\noindent\textbf{解析：} 一级页号 $10$ 位，页目录共 $2^{10}=1024$ 项，每项 $4$ B，占 $4$ KiB。两个虚页的一级页号分别是 $0$ 和 $255$（第二个虚页的二级页号为 $1023$），因此需要两张二级页表；每张二级页表有 $2^{10}$ 项，同样占 $4$ KiB。未被使用的一级页号不分配二级页表，所以总占用为
\[
4\text{ KiB}+2\times4\text{ KiB}=12\text{ KiB}.
\]
选项 (D) 是一级页表方案的大小：$2^{20}\times4\text{ B}=4$ MiB，它没有利用多级页表按需分配的特点。\boxed{\textbf{答案：}(C)}\else \answerarea \fi

\item \textbf{（单项选择题）} 某系统采用改进型 CLOCK 算法，4 个页框按顺时针顺序存放页面 $P_0,P_1,P_2,P_3$，其（访问位 $A$，修改位 $M$）依次为 $(1,1)$、$(0,1)$、$(1,0)$、$(0,1)$，时钟指针当前指向 $P_0$。按“第一轮寻找 $(0,0)$ 且不修改标志位，第二轮寻找 $(0,1)$ 并把扫过的 $A$ 清 $0$”的规则，下一个被换出的页面是
\begin{center}
\begin{tabular}{ll}
(A) $P_0$ & (B) $P_1$ \\[4pt]
(C) $P_2$ & (D) $P_3$
\end{tabular}
\end{center}

\ifshowsolutions\par\noindent\textbf{解析：} 四个页面中没有 $(A,M)=(0,0)$，第一轮扫描不选任何页。第二轮从 $P_0$ 开始：$P_0$ 为 $(1,1)$，把 $A$ 清 $0$ 后指针前移；$P_1$ 为 $(0,1)$，满足条件，立即选为牺牲页。因此被换出的是 $P_1$。若把第一轮误当成“边扫边清 $A$”，会在 $P_0$ 变 $(0,1)$ 时提前选 $P_0$，这是常见错误。\boxed{\textbf{答案：}(B)}\else \answerarea \fi

\item \textbf{（真题风格·综合题）} 某请求分页系统采用固定分配、局部置换，进程的页面访问序列为
\[
1,2,3,4,2,1,5,6,2,1,2,3,7,6,3,2,1,2,3,6,
\]
共 $20$ 次访问。规定：初始页框全空；缺页次数按“被访问页不在内存”计数，首次装入也算缺页；命中不改变 FIFO 队列的次序。
\begin{enumerate}
\item[(1)] 分配 $3$ 个页框，分别用 FIFO、LRU、OPT 求缺页次数与缺页率。
\item[(2)] 将页框数增加到 $4$，求 FIFO 与 LRU 的缺页次数；与 (1) 比较，该序列是否出现 Belady 异常？
\item[(3)] 另取访问序列 $1,2,3,4,1,2,5,1,2,3,4,5$，用 FIFO 分别计算 $3$ 个和 $4$ 个页框时的缺页次数，并解释 Belady 异常的成因。
\end{enumerate}

\ifshowsolutions\par\noindent\textbf{解析：}
(1) 逐次模拟结果如下。表中按页框槽位列出内容，带 $*$ 的是本次缺页新调入的页；OPT 在置换时选择“下次访问最晚”的页，若某页不再被访问则优先淘汰。

\begin{center}
\small
\begin{tabular}{|c|c|c|c|c|}
\hline
次序 & 访问页 & FIFO & LRU & OPT \\
\hline
1 & 1 & $1^*$ & $1^*$ & $1^*$ \\
2 & 2 & $1\,|\,2^*$ & $1\,|\,2^*$ & $1\,|\,2^*$ \\
3 & 3 & $1\,|\,2\,|\,3^*$ & $1\,|\,2\,|\,3^*$ & $1\,|\,2\,|\,3^*$ \\
4 & 4 & $4^*\,|\,2\,|\,3$ & $4^*\,|\,2\,|\,3$ & $1\,|\,2\,|\,4^*$ \\
5 & 2 & $4\,|\,2\,|\,3$ & $4\,|\,2\,|\,3$ & $1\,|\,2\,|\,4$ \\
6 & 1 & $4\,|\,1^*\,|\,3$ & $4\,|\,2\,|\,1^*$ & $1\,|\,2\,|\,4$ \\
7 & 5 & $4\,|\,1\,|\,5^*$ & $5^*\,|\,2\,|\,1$ & $1\,|\,2\,|\,5^*$ \\
8 & 6 & $6^*\,|\,1\,|\,5$ & $5\,|\,6^*\,|\,1$ & $1\,|\,2\,|\,6^*$ \\
9 & 2 & $6\,|\,2^*\,|\,5$ & $5\,|\,6\,|\,2^*$ & $1\,|\,2\,|\,6$ \\
10 & 1 & $6\,|\,2\,|\,1^*$ & $1^*\,|\,6\,|\,2$ & $1\,|\,2\,|\,6$ \\
11 & 2 & $6\,|\,2\,|\,1$ & $1\,|\,6\,|\,2$ & $1\,|\,2\,|\,6$ \\
12 & 3 & $3^*\,|\,2\,|\,1$ & $1\,|\,3^*\,|\,2$ & $3^*\,|\,2\,|\,6$ \\
13 & 7 & $3\,|\,7^*\,|\,1$ & $7^*\,|\,3\,|\,2$ & $3\,|\,7^*\,|\,6$ \\
14 & 6 & $3\,|\,7\,|\,6^*$ & $7\,|\,3\,|\,6^*$ & $3\,|\,7\,|\,6$ \\
15 & 3 & $3\,|\,7\,|\,6$ & $7\,|\,3\,|\,6$ & $3\,|\,7\,|\,6$ \\
16 & 2 & $2^*\,|\,7\,|\,6$ & $2^*\,|\,3\,|\,6$ & $3\,|\,2^*\,|\,6$ \\
17 & 1 & $2\,|\,1^*\,|\,6$ & $2\,|\,3\,|\,1^*$ & $3\,|\,2\,|\,1^*$ \\
18 & 2 & $2\,|\,1\,|\,6$ & $2\,|\,3\,|\,1$ & $3\,|\,2\,|\,1$ \\
19 & 3 & $2\,|\,1\,|\,3^*$ & $2\,|\,3\,|\,1$ & $3\,|\,2\,|\,1$ \\
20 & 6 & $6^*\,|\,1\,|\,3$ & $2\,|\,3\,|\,6^*$ & $6^*\,|\,2\,|\,1$ \\
\hline
\end{tabular}
\end{center}

FIFO 缺页发生在第 $1,2,3,4,6,7,8,9,10,12,13,14,16,17,19,20$ 次，共 $16$ 次，缺页率 $16/20=80\%$。LRU 缺页发生在第 $1,2,3,4,6,7,8,9,10,12,13,14,16,17,20$ 次，共 $15$ 次，缺页率 $75\%$。OPT 缺页发生在第 $1,2,3,4,7,8,12,13,16,17,20$ 次，共 $11$ 次，缺页率 $55\%$。

(2) 页框数改为 $4$ 时，FIFO 缺页 $14$ 次（$70\%$），LRU 缺页 $10$ 次（$50\%$）。两者都随页框增加而减少，因此\textbf{该序列没有出现} Belady 异常。Belady 异常并非 FIFO 的必然结果，只说明 FIFO 不具备保证缺页数单调不增的性质。

(3) 对序列 $1,2,3,4,1,2,5,1,2,3,4,5$：FIFO 在 $3$ 个页框时缺页 $9$ 次，在 $4$ 个页框时缺页 $10$ 次，页框增多反而多缺页，这就是 Belady 异常。原因是 FIFO 只按进入内存的先后淘汰，不看页面近期是否活跃；增加页框会改变淘汰次序，可能把即将访问的页提前换出。LRU 与 OPT 属于栈算法：$n$ 个页框的驻留集始终包含于 $n+1$ 个页框的驻留集，所以缺页数不会随页框增加而上升。本序列用 LRU 时 $3$ 个页框缺页 $10$ 次、$4$ 个页框缺页 $8$ 次，正好印证这一点。

\boxed{\textbf{答案：}(1)\ 16,\ 15,\ 11;\quad (2)\ 14,\ 10,\ \text{无异常};\quad (3)\ 9,\ 10,\ \text{有异常}}\else \answerarea[12] \fi

\item \textbf{（真题风格·综合题）} 某 32 位系统按字节编址，页面大小 $4$ KiB，采用二级页表，逻辑地址按“一级页号 $10$ 位、二级页号 $10$ 位、页内偏移 $12$ 位”划分，页表项与页目录项均为 $4$ B。页目录基址寄存器内容为 \texttt{0x00100000}。进程 $P$ 的映射如下：
\begin{itemize}
\item 一级页号 $1$ 的页目录项指向二级页表 \texttt{0x00200000}；
\item 在 \texttt{0x00200000} 这张二级页表中，二级页号 $3$ 对应页框号 \texttt{0x2A}，二级页号 $5$ 对应页框号 \texttt{0x37}；
\item 一级页号 $2$ 的页目录项指向的二级页表中，二级页号 $3$ 对应页框号 \texttt{0x11}。
\end{itemize}
\begin{enumerate}
\item[(1)] 访问逻辑地址 \texttt{0x00403ABC}，求一级页号、二级页号与页内偏移。
\item[(2)] 求读取该地址对应二级页表项时所用的物理地址，以及最终访问数据的物理地址。
\item[(3)] TLB 初始为空，依次访问 \texttt{0x00403ABC}、\texttt{0x00403AB8}、\texttt{0x00405ABC}、\texttt{0x00803ABC}，写出每次的 TLB 命中情况；并给出第一次访问在 TLB 缺失、页表命中时的访存次数。
\item[(4)] 若第 (1) 问地址的页表项有效位为 $0$，而该页合法且已分配在外存，说明缺页处理的主要步骤。再说明：若只给该进程 $1$ 个页框，而它循环访问 \texttt{0x00403ABC} 与 \texttt{0x00405ABC}，会出现什么现象，应如何缓解。
\end{enumerate}

\ifshowsolutions\par\noindent\textbf{解析：}
(1) 按位拆分：
\[
\texttt{0x00403ABC}=0000\,0000\,0100\,0000\,0011\,1010\,1011\,1100_2 .
\]
高 $10$ 位为 $0000000001_2=1$，接着 $10$ 位为 $0000000011_2=3$，低 $12$ 位为 \texttt{0xABC}。即一级页号 $1$、二级页号 $3$、页内偏移 \texttt{0xABC}（十进制 $2748$）。

(2) 页目录项地址为
\[
\texttt{0x00100000}+1\times4=\texttt{0x00100004},
\]
其内容给出二级页表基址 \texttt{0x00200000}。二级页表项地址为
\[
\texttt{0x00200000}+3\times4=\texttt{0x0020000C},
\]
其中页框号为 \texttt{0x2A}。物理地址由页框号与页内偏移拼接：
\[
\texttt{0x2A}\times2^{12}+\texttt{0xABC}
=\texttt{0x2A000}+\texttt{0xABC}=\texttt{0x2AABC}.
\]
同理，\texttt{0x00403AB8} 在同一页内，物理地址为 \texttt{0x2AAB8}。

(3) 判断是否命中只看虚页号（一级页号与二级页号合起来）：
\begin{enumerate}[label=(\arabic*)]
\item \texttt{0x00403ABC}：虚页 \texttt{0x403}，TLB 缺失，查页表后装入；
\item \texttt{0x00403AB8}：虚页同为 \texttt{0x403}，TLB 命中；
\item \texttt{0x00405ABC}：虚页 \texttt{0x405}（一级页号 $1$、二级页号 $5$），TLB 缺失；
\item \texttt{0x00803ABC}：虚页 \texttt{0x803}（一级页号 $2$、二级页号 $3$），TLB 缺失。
\end{enumerate}
第一次访问时 TLB 缺失，需读页目录项、读二级页表项、再取数据，共 $3$ 次内存访问；后续同页访问若 TLB 命中，只需 $1$ 次。

(4) 缺页处理的主要步骤：保存现场并确认该虚地址属于进程的合法地址空间；若页表项无效但页在外存，则分配空闲页框（无空闲时按置换算法选牺牲页，脏页先写回）；启动磁盘 I/O 把该页读入页框；更新页表项的有效位与页框号，必要时清空或更新相关 TLB 项；恢复现场并\textbf{重新执行}触发缺页的指令。

只给 $1$ 个页框时，两个地址分属虚页 \texttt{0x403} 与 \texttt{0x405}，程序每访问另一个页面都会把它们互相换出，形成“访问一次、缺页一次”的循环。此时进程几乎全部时间都在等待磁盘 I/O，CPU 利用率很低，属于抖动。页框数至少应达到该循环的工作集大小 $2$；若系统整体页框不足，则应减少并发进程数，或按工作集模型、缺页频率策略为进程增配页框。

\par\noindent\boxed{\textbf{答案：}(1)\ 1,3,\texttt{0xABC};\ (2)\ \texttt{0x0020000C},\texttt{0x2AABC};\ (3)\ 3\ \text{次};\ (4)\ \text{见解析}}\else \answerarea[14] \fi

\end{enumerate}
